\documentclass[verbale=interno]{gpt}

\titolo{Contatti con i proponenti e assegnazione dei ruoli}
\data{2026-10-19}
\orario{21:00}{22:30}
\luogo{Discord}
\redattore{guglielmo}
\verificatore{corro}
\componenti{biasini,cescobolla,corro,guglielmo,mazzurana,spagnoli}
\assenti{}

\modifica{1.0.0}{2026-10-19}{guglielmo}{corro}{Verifica e approvazione}
\modifica{0.1.0}{2026-10-19}{guglielmo}{--}{Prima stesura}

\begin{document}

\section{Ordine del giorno}
\begin{enumerate}
  \item Contatti con i proponenti dei capitolati C0, C1 e C2.
  \item Comunicazione dei dati del gruppo al committente.
  \item Creazione dell'organizzazione GitHub del gruppo.
  \item Assegnazione dei ruoli.
  \item Calendario delle riunioni.
  \item Formazione e prossimi passi.
\end{enumerate}

\section{Svolgimento}

\subsection{Contatti con i proponenti}
Durante l'incontro il gruppo ha avviato i contatti con i proponenti dei
tre capitolati di interesse:
\begin{itemize}
  \item \textbf{C2 -- NEXUM Concierge}: inviata un'email a Eggon S.r.l.;
  \item \textbf{C0 -- AIRA}: inviata un'email a Sanmarco Informatica S.p.A.;
  \item \textbf{C1 -- AURA}: compilato il form di Vimar S.p.A. per la
    richiesta di un incontro, indicando come preferenza giovedì
    22 ottobre 2026 dalle 16:00 alle 17:00. Il riferimento per la call è
    \url{https://aura.vimaraei.cloud/}.
\end{itemize}
Il gruppo è in attesa delle risposte dei proponenti e della conferma
della data da parte di Vimar.

\subsection{Comunicazione al committente}
È stata inviata un'email al Prof.~Tullio Vardanega con il nome del
gruppo, il recapito del riflettore di posta elettronica e il logo.

\subsection{Organizzazione GitHub}
È stata creata l'organizzazione GitHub del gruppo,
\textbf{Git Push Team} (\url{https://github.com/Git-Push-Team}). La
struttura dei repository non è ancora stata definita: gli
Amministratori si informeranno su come organizzarli.

\subsection{Assegnazione dei ruoli}
Il gruppo ha assegnato i ruoli di avvio, validi fino all'aggiudicazione
dei capitolati del 3 novembre 2026:
\begin{itemize}
  \item \textbf{Responsabile}: \nome{spagnoli};
  \item \textbf{Amministratori}: \nomi{guglielmo,mazzurana,biasini};
  \item \textbf{Verificatori}: \nomi{cescobolla,corro}.
\end{itemize}

\subsection{Calendario delle riunioni}
Il gruppo ha fissato una riunione ordinaria settimanale il lunedì alle
21:00, con la possibilità di variazioni da concordare in corso d'opera.
È inoltre prevista una riunione straordinaria in settimana, con data da
definire.

\subsection{Formazione e prossimi passi}
Ogni membro si impegna a formarsi sull'uso di GitHub prima del prossimo
incontro. Per facilitare questo lavoro, \nome{guglielmo} e
\nome{spagnoli} terranno una call straordinaria martedì 20 ottobre per
avviare l'organizzazione dei repository e preparare una guida a GitHub
per il resto del gruppo.

\section{Decisioni}
\begin{decisioni}
  \decisione{I ruoli di avvio, validi fino al 3 novembre 2026, sono:
    Responsabile \nome{spagnoli}; Amministratori
    \nomi{guglielmo,mazzurana,biasini}; Verificatori
    \nomi{cescobolla,corro}.}
  \decisione{La riunione ordinaria del gruppo si tiene ogni lunedì alle
    21:00, salvo variazioni concordate in corso d'opera.}
  \decisione{L'organizzazione GitHub del gruppo è Git Push Team.}
  \decisione{Ogni membro si forma sull'uso di GitHub prima del prossimo
    incontro.}
\end{decisioni}

\section{Azioni da intraprendere}
\begin{azioni}
  \azione{Call straordinaria per i primi passi nell'organizzazione dei
    repository e preparazione di una guida a GitHub per il
    gruppo.}{guglielmo,spagnoli}{20/10/2026}
  \azione{Informarsi su come organizzare i repository del
    gruppo.}{guglielmo,mazzurana,biasini}{26/10/2026}
  \azione{Formarsi sull'uso di GitHub.}{tutti}{26/10/2026}
  \azione{Concordare data e ora della riunione straordinaria della
    settimana.}{tutti}{25/10/2026}
\end{azioni}

\section{Prossima riunione}
Riunione straordinaria in settimana, con data da definire. La prossima
riunione ordinaria è lunedì 26 ottobre 2026 alle 21:00 su Discord.

\end{document}
